\subsection{Study canal and input data}

The study reach is the 0--456 main canal of the Shijin irrigation canal
prototype. The hydraulic model uses three raw input files. The file
\texttt{input.txt} provides node identifiers, planar coordinates, bed elevation
and channel style. The file \texttt{lineParam.txt} provides the cross-sectional
parameters associated with each channel style, including design depth, bottom
width, side-wall angle and Manning roughness. The file \texttt{neighborId.txt}
defines the topological connectivity between nodes. The numerical scenario
parameters, including time step, start and stop times, ramp duration, target
head flow and safety depth ratio, are read from
\texttt{data/configuration.json}.

The main canal is represented by the ordered nodes from 0 to 456. The total
main-canal length is 19941.40 m, with 457 computational nodes and 456
inter-node reaches. The finite-volume cell length is computed from the actual
node spacing. For interior cells,
\begin{equation}
    \Delta x_i = \frac{x_{i+1} - x_{i-1}}{2},
\end{equation}
whereas the two boundary cells use one-sided lengths. In the current data set,
the finite-volume cell length ranges from 34.25 m to 52.94 m, with a mean value
of 43.73 m. Seven diversion junctions are considered on the main canal:
71, 89, 150, 194, 287, 349 and 383. Their connected branch chains contain
12, 12, 30, 34, 11, 23 and 25 nodes, respectively.

\subsection{Cross-sectional hydraulic relations}

Each node is assigned a trapezoidal cross section from its channel style. For
water depth \(h\), bottom width \(b\), and side slope \(z\) expressed as
horizontal:vertical, the area, top width, wetted perimeter and hydraulic radius
are
\begin{equation}
    A(h) = bh + zh^2,
\end{equation}
\begin{equation}
    B(h) = b + 2zh,
\end{equation}
\begin{equation}
    P(h) = b + 2h\sqrt{1+z^2},
\end{equation}
\begin{equation}
    R(h) = \frac{A(h)}{P(h)}.
\end{equation}
The pressure integral used in the Saint-Venant momentum flux is
\begin{equation}
    I_1(h) = \frac{1}{2}bh^2 + \frac{1}{3}zh^3.
\end{equation}
Area and depth are converted analytically. Given area \(A\), the corresponding
depth is
\begin{equation}
    h(A) =
    \frac{-b + \sqrt{b^2 + 4zA}}{2z},
\end{equation}
with the rectangular limit \(h=A/b\) used when \(z\) is close to zero.

\subsection{Saint-Venant governing equations}

The forward hydraulic solver uses the one-dimensional Saint-Venant equations in
conservative form. The conserved state vector and flux vector are
\begin{equation}
    \mathbf{U} =
    \left[
    \begin{array}{c}
        A \\
        Q
    \end{array}
    \right],
    \qquad
    \mathbf{F}(\mathbf{U}) =
    \left[
    \begin{array}{c}
        Q \\
        Q^2/A + g I_1
    \end{array}
    \right],
\end{equation}
where \(Q\) is discharge and \(g=9.81\ {\rm m\,s^{-2}}\). The governing equation
for each channel segment is
\begin{equation}
    \frac{\partial \mathbf{U}}{\partial t}
    +
    \frac{\partial \mathbf{F}(\mathbf{U})}{\partial x}
    =
    \mathbf{S}_b + \mathbf{S}_f + \mathbf{S}_J,
\end{equation}
where \(\mathbf{S}_b\) is the bed-slope source term, \(\mathbf{S}_f\) is the
Manning friction source term, and \(\mathbf{S}_J\) is the junction exchange term
between the main canal and a branch canal.

The local bed-slope source term is
\begin{equation}
    \mathbf{S}_{b,i} =
    \left[
    \begin{array}{c}
        0 \\
        g A_i S_{0,i}
    \end{array}
    \right].
\end{equation}
For interior cells, \(S_{0,i}\) is computed from neighboring bed elevations:
\begin{equation}
    S_{0,i} =
    \frac{Z_{b,i-1} - Z_{b,i+1}}{x_{i+1}-x_{i-1}},
    \qquad 1 \le i \le N-1.
\end{equation}
At the two channel boundaries, one-sided bed slopes are used. A minimum slope of
\(5.0\times 10^{-5}\) is imposed to avoid unrealistically small Manning
capacity in nearly flat local segments.

\subsection{Fully implicit network discretization}

The current dispatch model uses a fully implicit graph-network discretization
instead of the earlier explicit HLL finite-volume update. Each computational
node has water level \(H_i=Z_{b,i}+h_i\) as an unknown, and each directed link
\(e=(i,j)\) has discharge \(Q_e\) as an unknown. The main canal, branch canals
and diversion connectors are therefore solved in a single network system.

For node \(i\), the implicit continuity equation is
\begin{equation}
    \frac{C_i^{(r)}}{\Delta t}
    \left(
        H_i^{n+1}-H_i^n
    \right)
    +
    \sum_{e\in{\cal O}(i)} Q_e^{n+1}
    -
    \sum_{e\in{\cal I}(i)} Q_e^{n+1}
    +
    Q_{{\rm out},i}^{n+1}
    -
    Q_{{\rm in},i}^{n+1}
    =0,
\end{equation}
where \({\cal O}(i)\) and \({\cal I}(i)\) are outgoing and incoming links,
respectively. \(C_i^{(r)}=B_i^{(r)}\Delta x_i\) is the linearized storage
coefficient at Picard iteration \(r\). The upstream boundary enters
\(Q_{{\rm in},i}^{n+1}\), while downstream open boundaries are represented by a
linearized Manning rating \(Q_{{\rm out},i}^{n+1}\).

For link \(e=(i,j)\), the implicit momentum relation is written in water-level
form as
\begin{equation}
    \frac{Q_e^{n+1}-Q_e^n}{\Delta t}
    +
    g A_e^{(r)}
    \left[
        \frac{H_j^{n+1}-H_i^{n+1}}{L_e}
        +
        \frac{
            n_e^2 |Q_e^n| Q_e^{n+1}
        }{
            \left(A_e^{(r)}\right)^2
            \left(R_e^{(r)}\right)^{4/3}
        }
    \right]
    =0 .
\end{equation}
This is the implicit local-inertial Saint-Venant link equation with Manning
friction linearized by a reference discharge based on the previous discharge;
a small normal-flow reference is used during cold starts to avoid a zero-friction
first step. Solving it for \(Q_e^{n+1}\) gives
\begin{equation}
    Q_e^{n+1}
    =
    q_e^{(r)}
    +
    K_e^{(r)}
    \left(
        H_i^{n+1}-H_j^{n+1}
    \right),
\end{equation}
with
\begin{equation}
    K_e^{(r)}
    =
    \frac{
        \Delta t\,g\,A_e^{(r)}/L_e
    }{
        1+
        \Delta t\,g\,n_e^2 |Q_e^n|
        /\left[
            A_e^{(r)}
            \left(R_e^{(r)}\right)^{4/3}
        \right]
    },
\end{equation}
\begin{equation}
    q_e^{(r)}
    =
    \frac{
        Q_e^n
    }{
        1+
        \Delta t\,g\,n_e^2 |Q_e^n|
        /\left[
            A_e^{(r)}
            \left(R_e^{(r)}\right)^{4/3}
        \right]
    } .
\end{equation}
Substituting the link expression into all node-continuity equations produces a
symmetric positive graph-Laplacian system for \(H^{n+1}\):
\begin{equation}
    \left[
        {\bf D}_C
        +
        {\bf D}_{\rm out}
        +
        {\bf L}({\bf K})
    \right]
    {\bf H}^{n+1}
    =
    {\bf b},
\end{equation}
where \({\bf D}_C\) is the storage diagonal matrix,
\({\bf D}_{\rm out}\) is the downstream rating derivative, and
\({\bf L}({\bf K})\) is the weighted graph Laplacian formed from the link
conductances \(K_e\). Because the current canal topology is a tree, the system
is solved by direct leaf-to-root elimination and root-to-leaf back substitution.
Picard iterations update the section area, hydraulic radius and storage
coefficients within each time step.

\subsection{Boundary conditions and simulation settings}

The upstream boundary is prescribed as a smooth inflow hydrograph that ramps
from 0 to 80 \({\rm m^3\,s^{-1}}\) during the first 0.5 h:
\begin{equation}
    Q_{\rm in}(t) =
    \left\{
    \begin{array}{ll}
        80 \times
        \frac{1-\cos(\pi t/T_r)}{2},
        & 0 < t < T_r, \\
        80, & t \ge T_r,
    \end{array}
    \right.
\end{equation}
where \(T_r=0.5\) h. The upstream inflow is inserted directly into the
upstream-node continuity equation. Downstream open boundaries use the local
Manning rating curve linearized at the current Picard iterate. The current
configuration uses a spatial resampling interval of 1 m, a computational time
step of \(\Delta t=600\) s, and a total simulation duration of 14 h.

\subsection{Main-branch junction coupling}

The earlier source-term-only representation of diversions was replaced by
network junction coupling. Each diversion \(k\) is represented as a directed
connector link from the main-canal node to the first branch node. Its flow is
therefore included in the same node-continuity equations and is driven by the
implicit water-level difference \(H_m^{n+1}-H_b^{n+1}\) through the link
momentum relation above. This treatment satisfies mass conservation at the
main-channel node and the branch inlet simultaneously, and it allows branch
backwater to reduce the diversion discharge through the coupled water levels.

The safe branch capacity is estimated at 90\% of the branch design depth:
\begin{equation}
    Q_{{\rm safe},k}
    =
    \frac{1}{n_b}
    A_b(0.9D_b)
    R_b(0.9D_b)^{2/3}
    S_{b,k}^{1/2}.
\end{equation}
The water-depth activation factor at the diversion gate is
\begin{equation}
    \phi(h) =
    \left\{
    \begin{array}{ll}
        0, & h \le 0.2D, \\
        \sqrt{\frac{h-0.2D}{0.6D-0.2D}}, & 0.2D < h < 0.6D, \\
        1, & h \ge 0.6D.
    \end{array}
    \right.
\end{equation}
The actual diversion flow is constrained by
\begin{equation}
    0
    \le
    Q_{{\rm div},k}^{n+1}
    \le
    Q_{{\rm lim},k}
    =
    \phi(h_{m,k})
    \min
    \left(
        Q_{{\rm max},k},
        Q_{{\rm safe},k},
        \frac{W_{{\rm rem},k}}{\Delta t}
    \right),
\end{equation}
Here \(Q_{{\rm max},k}\) is the specified outlet capacity, \(W_{{\rm rem},k}\)
is the remaining demand, and \(Q_{{\rm lim},k}\) is the operating upper bound
used by the implicit active set. If the unconstrained connector flow is
negative, the connector is closed. If it exceeds \(Q_{{\rm lim},k}\), it is
treated as a fixed-flow constrained link. Otherwise it remains an unconstrained
implicit momentum link. The active set is repeated within the time step until
the connector states no longer change.

\subsection{Dispatch assumptions}

The diversion outlets at nodes 71, 89, 150, 194, 287, 349 and 383 are assigned
maximum capacities of 20, 20, 5, 5, 12, 5 and 5 \({\rm m^3\,s^{-1}}\),
respectively. The corresponding demand volumes are 250000, 300000, 60000,
80000, 160000, 50000 and 40000 \({\rm m^3}\). The cumulative supplied volume is
\begin{equation}
    W_k(t_m) =
    \sum_{j=1}^{m}
    Q_{{\rm div},k}(t_j)\Delta t.
\end{equation}
Once \(W_k\) reaches the prescribed demand, the outlet is closed and its
diversion flow is set to zero for the remainder of the simulation.

\subsection{NumPy acceleration}

The numerical kernel is implemented in Python with NumPy arrays. The 1 m
resampled main canal and branch chains are stored as one graph. Link endpoints,
link lengths, section parameters and connector ownership are all stored as
index arrays. The implicit graph matrix has the same tree topology as the canal
network, so the linear solve is performed by one post-order elimination sweep
followed by one pre-order back-substitution sweep. The cost per Picard or active
set iteration scales linearly with the number of network links. For the current
1 m grid, the main canal contains about 20200 computational nodes and the routed
branches contain about 6200 nodes.

\subsection{Stability assessment and postprocessing}

The previous explicit HLL solver had to satisfy the Courant condition
\begin{equation}
    {\rm CFL}
    =
    \max_i
    \frac{(|u_i|+c_i)\Delta t}{\Delta x_i}.
\end{equation}
For \(\Delta x=1\) m and \(\Delta t=600\) s, this CFL number is several
thousand under the 80 \({\rm m^3\,s^{-1}}\) flow scale, so the explicit solver
is not suitable. The fully implicit network solver removes this explicit CFL
restriction for numerical stability by solving the coupled water levels and
link discharges at the new time level. The time step is therefore controlled
primarily by nonlinear convergence, boundary-condition resolution and physical
accuracy rather than by the explicit wave-propagation CFL limit.

Postprocessing is separated from simulation. The simulator module runs the
network-coupled hydraulic and dispatch model only. The postprocess module
generates diversion hydrographs, cumulative supplied volumes, key-node water
depth time series, and maximum-depth envelope diagnostics. The maximum-depth
envelope reuses the same network-coupled dispatch result and reports the largest
water depth reached at every main-canal node during the simulation.
